\section{Datasets}
\label{sec:datasets}

We use published datasets only; no data is self-collected. Table~%
\ref{tab:datasets} summarises the datasets and their programming-specific
signals.

\begin{table}[t]
\caption{Datasets used in this study.}
\label{tab:datasets}
\centering
\small
\begin{tabular}{@{}llll@{}}
\toprule
\textbf{Dataset} & \textbf{Domain} & \textbf{KC tags} & \textbf{Code} \\
\midrule
CodeWorkout (CSEDM'21) & Java CS1     & Yes & Yes \\
FalconCode            & Python intro & Yes & Yes \\
Code.org (optional)   & Blocks       & Partial & Partial \\
ASSISTments / Junyi   & Mathematics  & Yes & No  \\
\bottomrule
\end{tabular}
\end{table}

\subsection{CodeWorkout (CSEDM 2021)}
Collected from a CS1 course over two semesters (Spring and Fall~2019),
with 329 and 490 students respectively completing the course; each
semester contains 65{,}000+ code submissions across 50 Java coding
problems, each requiring roughly 10--26 lines of
code~\cite{csedm2021}. It is distributed in the ProgSnap2 format and
includes knowledge-component tags, making it our primary programming
dataset.

\subsection{FalconCode}
A multi-year dataset of introductory Python programming containing over
1.5 million code submissions from more than 2{,}000 undergraduate
students across 800+ assignments, with prompts, test cases, and
per-problem skills~\cite{falconcode2023}. It serves as our second
programming dataset for cross-dataset generalisation.

\subsection{Mathematics Contrast}
For the cross-domain comparison we use a standard mathematics KT dataset
(ASSISTments or Junyi) as preprocessed by pyKT~\cite{liu2022pykt}, which
provides KC tags but no code signal.

\subsection{Preprocessing and Splits}
All datasets are standardised into a common interaction format
(\texttt{student\_id}, \texttt{problem\_id}, \texttt{kc\_id},
\texttt{correct}, \texttt{code}). We adopt student-level
cross-validation (5-fold) so that no student appears in both training and
test folds, and we guard against the known label-leakage pitfall in KT
evaluation. The exact per-dataset counts after filtering (retained
students, interactions, KCs, and the sequence-length cap) are reported in
Table~\ref{tab:datastats} once preprocessing is finalised; the
dataset-level totals above are taken from the original data
releases~\cite{csedm2021,falconcode2023}.

\begin{table}[t]
\caption{Dataset-level statistics from the original public
releases~\cite{csedm2021,falconcode2023}. Counts are approximate and
reflect the released data before any project-specific filtering.}
\label{tab:datastats}
\centering
\small
\begin{tabular}{@{}lrrr@{}}
\toprule
\textbf{Dataset} & \textbf{Students} & \textbf{Submissions}
                 & \textbf{Problems} \\
\midrule
CodeWorkout & $819$      & $130{,}000+$ & $50$ \\
FalconCode  & $2{,}000+$ & $1{,}500{,}000+$ & $800+$ \\
\bottomrule
\end{tabular}
\end{table}
